% mazzi: 30,31
\chapter{Il sistema nervoso centrale}

% ---------------------------------------------------------------------------
\section{Generalità}

\textbf{SNC}: rappresentato anatomicamente dal \textit{nevrasse}, suddivisibile
macroscopicamente in due parti:
\begin{itemize}
  \item \textbf{midollo spinale}: situato nel canale vertebrale;
  \item \textbf{encefalo}: localizzato nella cavità cranica.
\end{itemize}

Le principali strutture dell'encefalo sono, a loro volta, identificate in
tronco encefalico, cervelletto, diencefalo e telencefalo. Midollo spinale ed
encefalo sono in continuità anatomica a livello del forame magno, pur
presentando caratteristiche diverse di configurazione esterna e conformazione
interna.

\textbf{SNC}: circondato da 3 involucri o strati fibrovascolari, le
\textbf{meningi}; al suo interno presenta una serie di cavità comunicanti,
formate dal \textbf{canale centrale} nel midollo spinale e dai
\textbf{ventricoli encefalici} nell'encefalo. Il compartimento esterno del
nevrasse è rappresentato dallo \textbf{spazio subaracnoideo}.

\rubrica{Le regioni funzionali principali}
\begin{itemize}
  \item \textbf{talamo}: nucleo di integrazione e ritrasmissione;
  \item \textbf{subtalamo}: avvio e arresto del gesto motorio;
  \item \textbf{metatalamo}: trasmissione corticale di informazioni visive e
    uditive;
  \item \textbf{telencefalo}: percezione cosciente di stimoli esterni,
    comandi motori, aree associative (integrazione delle informazioni
    sensitive per l'interpretazione della realtà, aspetti emotivi, logici e
    linguistici);
  \item \textbf{epitalamo}: ritmi circadiani;
  \item \textbf{ipotalamo}: modulazione della vita vegetativa e controllo
    dell'ipofisi;
  \item \textbf{cervelletto}: coordinazione dei movimenti e postura;
  \item \textbf{tronco encefalico}: integrazione e ritrasmissione, riflessi
    del tronco;
  \item \textbf{midollo spinale}: riflessi spinali e ritrasmissione delle
    vie sensitive e motorie.
\end{itemize}

\rubrica{Sintesi per struttura}
\begin{itemize}
  \item \textbf{midollo spinale}: sinapsi con il sistema nervoso periferico,
    vie ascendenti (sensitive) e vie discendenti (motorie);
  \item \textbf{bulbo}: incrocio delle vie discendenti motorie (\textit{incrocio
    piramidale}), incrocio delle vie ascendenti sensitive, formazione
    reticolare, emergenza dei nervi cranici;
  \item \textbf{ponte}: peduncoli cerebellari, formazione reticolare,
    emergenza dei nervi cranici;
  \item \textbf{mesencefalo}: nuclei della base;
  \item \textbf{diencefalo}: ipofisi, ipotalamo, talamo, chiasma ottico;
  \item \textbf{telencefalo}: corteccia cerebrale;
  \item \textbf{cervelletto}: centro di integrazione dell'equilibrio.
\end{itemize}

% ---------------------------------------------------------------------------
\section{Cenni di sviluppo ed evoluzione}

Il cervello si sviluppa da un \textbf{tubo neurale} cavo, formato dal tessuto
dell'embrione nei primissimi stadi.
\begin{itemize}
  \item le cellule si moltiplicano rapidamente nella parte anteriore del
    tubo, che cresce diventando un cervello embrionale;
  \item le cellule appena prodotte diventano neuroni non ancora maturi;
  \item a circa 4 settimane di vita dell'embrione, migrano verso le loro
    destinazioni, sviluppando dendriti, assoni e formando milioni di
    connessioni sinaptiche;
  \item la moltiplicazione avviene senza un progetto definito: il cervello
    embrionale crea neuroni e sinapsi in eccesso, sono poi la competizione e
    le interazioni ambientali a determinare i circuiti funzionali;
  \item circa 1 neurone su 2 muore, non riuscendo a formare connessioni
    utili;
  \item gli assoni dei neuroni sopravvissuti vengono isolati dalle cellule
    gliali in un processo detto \textbf{mielinizzazione}, che aumenta
    velocità e qualità delle trasmissioni.
\end{itemize}
Lo sviluppo culmina nella prima adolescenza e si riduce nell'adulto.

\rubrica{Evoluzione}
Circa 1 miliardo di anni fa comparvero i primi organismi pluricellulari: le
loro cellule dovevano comunicare tra loro e svilupparono reti neurali, in una
sorta di protocervello ancora oggi presente in creature come le meduse.
Eventi geologici e climatici crearono nuovi ambienti, dando un'ulteriore
spinta all'evoluzione cerebrale con la comparsa di gruppi di neuroni
specializzati. L'uomo di Neanderthal possedeva un encefalo di dimensioni e
peso maggiori rispetto a quello dell'uomo moderno; l'\textit{Homo sapiens}
ha un encefalo di peso pari a circa 1600\,g nel maschio e 1450\,g nella
femmina.

% ---------------------------------------------------------------------------
\section{Le quattro principali regioni del cervello}

\begin{itemize}
  \item emisferi cerebrali;
  \item diencefalo;
  \item tronco encefalico: mesencefalo, ponte, midollo allungato (bulbo);
  \item cervelletto.
\end{itemize}

\rubrica{Il diencefalo}
\begin{itemize}
  \item circondato dagli emisferi cerebrali;
  \item composto da 3 strutture simmetriche: talamo, ipotalamo ed epitalamo
    (comprende l'ipofisi);
  \item confina con il $3^\circ$ ventricolo;
  \item principalmente costituito da materia grigia.
\end{itemize}

\rubrica{Il tronco encefalico}
Costituito da strutture complesse, svolge innumerevoli funzioni.
\begin{itemize}
  \item \textbf{mesencefalo}: formato da tetto e tegmento. Il tegmento
    controlla i movimenti dell'occhio ed è coinvolto nell'elaborazione degli
    stimoli dolorifici; contiene nuclei che elaborano informazioni visive e
    uditive e generano le risposte riflesse a questi stimoli;
  \item \textbf{ponte}: protuberanza del tronco cerebrale, costituisce parte
    del pavimento del $4^\circ$ ventricolo e posteriormente si unisce al
    cervelletto tramite il peduncolo cerebellare medio;
  \item \textbf{midollo allungato (bulbo)}: il midollo spinale si modifica
    progressivamente e diventa più complesso, articolandosi infine nel bulbo.
    A questo livello si incrociano le vie piramidali; contiene i centri per
    la regolazione delle funzioni viscerali (respirazione, pressione
    sanguigna).
\end{itemize}

\rubrica{Il mesencefalo in dettaglio}
\begin{itemize}
  \item \textbf{sostanza grigia}:
  \begin{itemize}
    \item \textit{tetto}: collicolo superiore (integra le informazioni
      visive con altri input sensitivi, inizia le risposte riflesse agli
      stimoli visivi), collicolo inferiore (trasmette le informazioni
      uditive ai nuclei genicolati mediali, inizia le risposte riflesse
      agli stimoli uditivi);
    \item \textit{pareti e pavimento}: nucleo rosso (controllo involontario
      del tono muscolare e della posizione degli arti), substantia nigra
      (regola l'attività dei nuclei della base), formazione reticolare
      (elaborazione automatica delle afferenze sensitive e delle efferenze
      motorie, può iniziare la risposta motoria involontaria a uno stimolo,
      opera il mantenimento dello stato di coscienza), altri nuclei/centri
      associati al III e al IV nervo cranico.
  \end{itemize}
  \item \textbf{sostanza bianca}: peduncoli cerebrali, che connettono la
    corteccia motoria primaria con i motoneuroni dell'encefalo e del midollo
    spinale e trasportano le informazioni sensitive ascendenti al talamo.
\end{itemize}

\rubrica{Funzioni generali del tronco}
\begin{itemize}
  \item produce comportamenti automatizzati necessari per la sopravvivenza;
  \item è via di passaggio per tutti i tratti di fibre tra cervello e
    midollo spinale;
  \item è coinvolto nell'innervazione della faccia e del capo;
  \item 10 delle 12 paia di nervi cranici partono da esso.
\end{itemize}

% ---------------------------------------------------------------------------
\section{Cenni storici}

Il \textit{Papiro Chirurgico di Edwin Smith}, risalente al XVII secolo a.C.,
contiene i primi riferimenti scritti relativi al cervello.

L'anatomista Galeno dissezionò numerosi cervelli di vari animali,
distinguendo l'encefalo, suddiviso dalle meningi, in ``cervello''
(\textit{encephalon}, o ``cervello anteriore'') e in ``cervello posteriore''
(\textit{enkranion} o \textit{epikranion}), cioè in cervello e cervelletto;
scoprì inoltre la presenza dei ventricoli.

% ---------------------------------------------------------------------------
\section{Termini direzionali del SNC}

\begin{itemize}
  \item \textbf{rostrale}: verso il naso (anteriore, verso l'alto);
  \item \textbf{caudale}: verso la coda (posteriore, verso il basso);
  \item \textbf{dorsale}: posteriore;
  \item \textbf{ventrale}: anteriore.
\end{itemize}

% ---------------------------------------------------------------------------
\section{Gli emisferi cerebrali e la corteccia}

I due emisferi sono separati da un profondo solco longitudinale (\textit{fessura
cerebrale longitudinale}, \textit{scissura interemisferica} o \textit{solco
sagittale}).

\rubrica{Corteccia cerebrale}
\begin{itemize}
  \item formata da neuroni, glia e fibre nervose senza mielina, con uno
    spessore di circa 2--4\,mm;
  \item gioca un ruolo centrale nelle funzioni mentali cognitive complesse
    come pensiero, consapevolezza, memoria, attenzione, linguaggio;
  \item in anatomia patologica assume un colore grigio, da cui il nome di
    \textit{sostanza grigia}.
\end{itemize}

\rubrica{I lobi e le loro funzioni}
\begin{itemize}
  \item \textbf{lobo frontale} (corteccia motoria primaria): controllo
    volontario della muscolatura scheletrica;
  \item \textbf{lobo parietale} (corteccia sensitiva primaria): percezione
    cosciente del tatto, della pressione, delle vibrazioni, del dolore,
    della temperatura e del gusto;
  \item \textbf{lobo occipitale} (corteccia visiva): percezione cosciente
    degli stimoli visivi;
  \item \textbf{lobo temporale} (corteccia uditiva e olfattiva): percezione
    cosciente degli stimoli uditivi e olfattivi;
  \item \textbf{tutti i lobi} (aree associative): integrazione ed
    elaborazione dei dati sensitivi, elaborazione e inizio delle attività
    motorie.
\end{itemize}
Oltre ai quattro lobi principali (frontale, parietale, temporale,
occipitale) si distingue anche un lobo limbico. Nella superficie laterale
dell'emisfero, in profondità rispetto al solco laterale, si trova inoltre
l'\textit{insula}, sede della corteccia gustativa.

\rubrica{Specializzazioni emisferiche}
\begin{itemize}
  \item \textbf{emisfero sinistro}: linguaggio, matematica, logica;
  \item \textbf{emisfero destro}: abilità spaziale, riconoscimento del
    viso, immaginazione visuale, musica; analizza le informazioni sensitive
    e mette il corpo in relazione con l'ambiente esterno. I centri
    dell'interpretazione presenti in questo emisfero permettono
    l'identificazione di oggetti familiari tramite tatto, odorato e gusto.
\end{itemize}

\rubrica{Sostanza bianca}
\begin{itemize}
  \item \textbf{fibre associative}: connettono le aree corticali nello
    stesso emisfero;
  \begin{itemize}
    \item \textit{fibre arcuate}: connettono le circonvoluzioni all'interno
      di un lobo;
    \item \textit{fascicoli longitudinali}: connettono il lobo frontale con
      gli altri lobi cerebrali;
  \end{itemize}
  \item \textbf{fibre commessurali} (commessura anteriore e corpo calloso):
    connettono i lobi corrispondenti dei due emisferi;
  \item \textbf{fibre di proiezione}: connettono la corteccia cerebrale al
    diencefalo, al tronco encefalico, al cervelletto e al midollo spinale.
\end{itemize}

\rubrica{Il cervello in cifre}
Il cervello ha forma simile a un cavolfiore, un peso di circa 1,36\,kg e una
consistenza paragonabile a quella del tofu. Le moderne tecnologie di
imaging (TC, PET, risonanza magnetica), insieme a genetica, chimica e
informatica, permettono di svelarne i particolari in modo sempre più
dettagliato.

% ---------------------------------------------------------------------------
\section{Le meningi encefaliche}

Sono la \textbf{dura madre}, l'\textbf{aracnoide} e la \textbf{pia madre}; di
queste, aracnoide e pia madre hanno la stessa struttura delle meningi
spinali, mentre la dura madre è costituita da due strati senza
interposizione dello spazio epidurale.

\rubrica{Dura madre}
Formata da un doppio strato di tessuto connettivo denso irregolare:
\begin{itemize}
  \item \textbf{strato periostale esterno}: equivalente al periostio delle
    ossa craniche;
  \item \textbf{strato meningeo interno}: solo questo strato continua nel
    midollo spinale.
\end{itemize}

\rubrica{Aracnoide e pia madre}
\begin{itemize}
  \item \textbf{aracnoide}: strato intermedio, costituito da epitelio
    pavimentoso semplice; con i suoi filamenti rimane saldo alla dura madre;
  \item \textbf{pia madre}: strato più esile e interno, a contatto indiretto
    (per la presenza di cellule gliali) con il cervello; costituita da un
    connettivo lasso, con fasci collagene ad andamento circolare.
\end{itemize}
Insieme, aracnoide e pia madre costituiscono la \textbf{leptomeninge}, che
funge da emuntorio con funzioni di controllo della produzione e del
riassorbimento del liquor.

\rubrica{Funzioni delle meningi}
\begin{itemize}
  \item proteggono l'encefalo e offrono supporto strutturale al passaggio
    delle arterie e delle vene encefaliche;
  \item offrono stabilità fisica e protezione dagli urti;
  \item aracnoide e pia madre, interponendosi tra i vasi e il tessuto
    nervoso, costituiscono la \textbf{barriera meningea}, che impedisce a
    sostanze tossiche, metaboliti e farmaci di penetrare dal sangue
    all'ambiente perineuronale, e che nutre il tessuto cerebrale.
\end{itemize}

\rubrica{Spazi intermeningei}
\begin{itemize}
  \item tra aracnoide e pia madre: \textbf{spazio subaracnoideo}, occupato
    dal liquido cerebrospinale (liquor);
  \item esternamente alla dura madre: \textbf{spazio epidurale}, che la
    separa dalla superficie ossea e contiene tessuto adiposo e plessi
    venosi.
\end{itemize}

\rubrica{Emorragia subdurale}
La rigida struttura dei seni venosi della dura madre, all'interno della
scatola cranica, comporta una serie di pericoli: un trauma violento (es.\
incidente stradale) a livello del cranio sposta bruscamente il contenuto
contro la dura madre, che è a contatto con l'osso. Le vene cerebrali vengono
compresse e possono andare incontro a rottura; spesso si verificano
lacerazioni a livello dello sbocco del seno sagittale superiore, con lenta
emorragia nel territorio compreso tra dura madre e aracnoide sottostante
(\textit{emorragia subdurale}). Questa si sviluppa spesso dopo giorni o
settimane e comprime in modo graduale le regioni cerebrali; il paziente
soffre di cefalea e stato confusionale, a causa dell'aumento della pressione
endocranica.

% ---------------------------------------------------------------------------
\section{Il sistema liquorale}

Le meningi, unitamente al liquido cerebrospinale, mantengono immerso e in
sospensione il sistema nervoso centrale, svolgendo una funzione protettiva.

\rubrica{Il liquor}
\begin{itemize}
  \item prodotto dai \textbf{plessi corioidei} situati nei ventricoli
    cerebrali, che lo filtrano dal sangue;
  \item percorso da trabecole fibrose che lo fissano alla pia madre;
  \item riassorbito dalle \textbf{villi aracnoidei} (o granulazioni
    aracnoidee di Pacchioni), che si trovano nei seni venosi cerebrali, in
    particolare nel seno sagittale, e lo immettono nel circolo venoso;
  \item lo spazio ventricolare e subaracnoideo fornisce lo spazio per la sua
    circolazione.
\end{itemize}
Bagna, isola, drena e nutre ogni parte del SNC, creando: un ambiente
ottimale per la riproduzione delle cellule della nevroglia e il
funzionamento delle cellule nervose; una protezione dai traumi esterni,
assorbendo e distribuendo le forze applicate sulla superficie dell'encefalo.

\rubrica{Circolazione del liquor}
Dopo aver attraversato le cavità ventricolari e il canale dell'ependima
midollare, il liquor esce nello spazio subaracnoideo attraverso l'apertura
mediana (forame di Magendie) e le aperture laterali (forami di Luschka) del
$4^\circ$ ventricolo. Risale poi negli spazi subaracnoidei, riempiendo le
varie cisterne, fino alle granulazioni aracnoidee del seno sagittale
superiore, dove viene riassorbito nel circolo venoso.

Diffonde inoltre passivamente attraverso la continuità delle membrane gliali
ed ependimali dei ventricoli, entrando nello spazio extracellulare e
unendosi al liquido extracellulare proveniente dal letto capillare.

\rubrica{Prelievo e caratteristiche}
\begin{itemize}
  \item si preleva tramite \textbf{puntura lombare}, effettuata tra $L3$ e
    $L4$ o tra $L4$ e $L5$ (a livello della cauda equina);
  \item viene riciclato circa 3 volte al giorno;
  \item il volume prodotto è di circa 400--500\,ml/die;
  \item si presenta di colore limpido.
\end{itemize}

% ---------------------------------------------------------------------------
\section{La barriera ematoencefalica}

L'acqua metabolica liquorale, i neurotrasmettitori e i metaboliti in essa
contenuti non sono in grado di superare le membrane gliali ed ependimali,
che costituiscono così una \textbf{barriera ematoencefalica} (BEE), garante
del mantenimento di un ambiente ideale e stabile per la popolazione
neuronale.

\textbf{BEE}: unità anatomico-funzionale realizzata dalle particolari
caratteristiche delle cellule endoteliali che compongono i vasi del SNC. Ha
funzione di protezione del tessuto cerebrale dagli elementi nocivi (chimici)
presenti nel sangue, pur permettendo il passaggio di sostanze necessarie
alle funzioni metaboliche.

Fu scoperta più di 100 anni fa: iniettando un colorante blu in circolo,
tutto il corpo diventava blu, ad eccezione del cervello. La BEE è
semipermeabile: alcune molecole passano, altre no (es.\ il glutammato
monosodico, coinvolto nella cosiddetta ``sindrome del ristorante cinese'').
Microonde, radiazioni, traumi e iperosmolarità possono danneggiare la
barriera.

% ---------------------------------------------------------------------------
\section{I ventricoli encefalici}

\begin{itemize}
  \item ripieni di liquido cerebrospinale;
  \item rivestiti di cellule ependimali;
  \item in continuità fra loro;
  \item in continuità con il canale centrale del midollo spinale.
\end{itemize}

L'encefalo presenta 4 cavità o ventricoli, al cui interno scorre il liquor
prodotto dai capillari dei plessi coroidei.
\begin{itemize}
  \item i \textbf{ventricoli laterali} sono i più grandi: uno per ciascun
    emisfero, comunicano con il $3^\circ$ ventricolo tramite il forame
    interventricolare (forame di Monro);
  \item il \textbf{$3^\circ$ ventricolo} è posto in posizione mediana, al di
    sotto del corpo calloso;
  \item l'\textbf{acquedotto del mesencefalo} (o acquedotto di Silvio)
    attraversa la porzione centrale del mesencefalo e si porta nel
    \textbf{$4^\circ$ ventricolo}, una cavità triangolare più piccola posta
    tra il ponte e il cervelletto, in continuità con il canale centrale del
    midollo spinale.
\end{itemize}

% ---------------------------------------------------------------------------
\section{La circolazione del sangue nel cervello}

Il cervello rappresenta il 2\% del peso corporeo totale, ma riceve il 15--20\% di
tutto il sangue in circolazione: il flusso deve restare sempre costante, perché i
neuroni sono cellule molto esigenti per il rifornimento di nutrienti ed energia.

\rubrica{Cosa porta e cosa rimuove il sangue}
\begin{itemize}
  \item il sangue \textbf{porta} al cervello: $O_2$, carboidrati, amminoacidi,
    grassi, ormoni, vitamine;
  \item il sangue \textbf{rimuove} dal cervello: biossido di carbonio, ammoniaca,
    lattato, ormoni.
\end{itemize}

\rubrica{Consumo metabolico}
\begin{itemize}
  \item l'encefalo è un organo ad alto consumo metabolico, con limitate risorse
    energetiche di glucosio e $O_2$: la vascolarizzazione arteriosa deve quindi
    essere costante e metabolismo-dipendente;
  \item al cervello arrivano circa 750\,cc di sangue al minuto (50\,cc/100\,g,
    contro i 100\,cc/100\,g del cuore);
  \item il cervello consuma circa 50\,cc di $O_2$/min (quasi il 15\% del totale)
    e 80\,mg/min di glucosio.
\end{itemize}

% ---------------------------------------------------------------------------
\section{Vascolarizzazione arteriosa dell'encefalo}

Alla base della scatola cranica è presente un sistema di anastomosi arteriose a
pieno canale, dato dalla confluenza di 3 arterie principali:
\begin{itemize}
  \item \textbf{l'arteria basilare}: formata dalla confluenza delle arterie
    vertebrali destra e sinistra (prime collaterali della succlavia);
  \item \textbf{le 2 arterie carotidi interne} (destra e sinistra).
\end{itemize}

\rubrica{Il poligono di Willis}
Può essere ricondotto idealmente a un eptagono, i cui lati sono:
\begin{itemize}
  \item anteriormente, le 2 arterie cerebrali anteriori (dx e sx), unite
    dall'arteria comunicante anteriore;
  \item posteriormente, le 2 arterie cerebrali posteriori (dx e sx);
  \item tra arteria cerebrale anteriore e posteriore di ogni lato, l'arteria
    comunicante posteriore.
\end{itemize}

\rubrica{Varianti anatomiche}
Il sistema arterioso può presentare varianti:
\begin{itemize}
  \item a livello di origine;
  \item di decorso (tortuosità);
  \item di calibro;
  \item del sifone carotideo.
\end{itemize}

\rubrica{Occlusione della circolazione cerebrale}
Quando la circolazione del sangue nel cervello è bloccata si parla di ictus,
trombosi, embolo o stenosi.

% ---------------------------------------------------------------------------
\section{I seni venosi}

I seni venosi sono canali a pareti rigide situati nello spessore della dura
madre, rivestiti da endotelio.

\rubrica{Caratteristiche}
\begin{itemize}
  \item volume più considerevole rispetto alle arterie;
  \item i rami decorrono di preferenza alla superficie libera delle
    circonvoluzioni;
  \item pareti molto sottili, prive di fibre muscolari;
  \item assenza di valvole;
  \item molteplicità di anastomosi;
  \item sono tributari, direttamente o indirettamente, delle vene giugulari
    interne;
  \item ricevono sangue dalle vene cerebrali.
\end{itemize}

% ---------------------------------------------------------------------------
\section{I nuclei della base}

Localizzati nel telencefalo basale, costituiscono nel loro insieme il
\textbf{corpo striato}, suddiviso filogeneticamente in:
\begin{itemize}
  \item \textbf{striato} o neo-striato: nucleo caudato e putamen;
  \item \textbf{pallido} o paleo-striato: \textit{globus pallidus} laterale e
    mediale.
\end{itemize}

Rappresentano un importante sistema di selezione, avvio e arresto del
movimento volontario.

\rubrica{Circuiti}
Motorio, oculomotorio, cognitivo, limbico.

\rubrica{Funzioni}
\begin{itemize}
  \item controllo inconscio e integrazione del tono della muscolatura
    scheletrica;
  \item coordinazione degli schemi motori acquisiti;
  \item elaborazione, integrazione e trasmissione delle informazioni dalla
    corteccia cerebrale verso il talamo.
\end{itemize}

\rubrica{Patologie dei nuclei della base}
Sindromi rigido-acinetiche, sindromi iper-acinetiche, disturbi del movimento.

% ---------------------------------------------------------------------------
\section{Il sistema limbico}

Comprende nuclei e fasci che si trovano lungo il confine tra cervello e
diencefalo: rappresenta un raggruppamento più funzionale che anatomico, legato
al comportamento dell'individuo.

\rubrica{Componenti}
\begin{itemize}
  \item \textbf{funzioni}: elaborazione dei ricordi, creazione degli stati
    emozionali, della condotta e del comportamento;
  \item \textbf{componenti cerebrali}:
    \begin{itemize}
      \item aree corticali -- lobo limbico (giro del cingolo, giro
        paraippocampale, giro dentato);
      \item nuclei -- ippocampo e corpo amigdaloideo;
      \item tratti -- fornice;
    \end{itemize}
  \item \textbf{componenti diencefaliche}:
    \begin{itemize}
      \item talamo -- gruppo nucleare anteriore;
      \item ipotalamo -- centri delle emozioni, dell'appetito (fame e sete) e
        del relativo comportamento;
    \end{itemize}
  \item \textbf{altre componenti}: formazione reticolare, rete di nuclei
    cerebrali interconnessi in tutto il tronco encefalico.
\end{itemize}

% ---------------------------------------------------------------------------
\section{Il talamo}

Rappresenta il punto di trasmissione finale per le informazioni sensitive e
coordina le attività motorie somatiche volontarie e involontarie. Contiene
nuclei che rappresentano centri di controllo delle informazioni sensitive e
motorie.

\begin{itemize}
  \item forma simile a un uovo: è il ``relè'' del cervello;
  \item struttura posta tra tronco encefalico ed emisferi cerebrali;
  \item quasi tutte le informazioni sensoriali passano da qui prima di essere
    inviate alla corteccia;
  \item quasi tutti i nuclei talamici sono connessi con la corteccia.
\end{itemize}

% ---------------------------------------------------------------------------
\section{L'ipotalamo}

Controlla numerose funzioni autonome, tra cui l'omeostasi e il controllo
ipofisario: è il centro di collegamento tra sistema nervoso e sistema
endocrino.

\rubrica{Funzioni}
\begin{enumerate}
  \item \textbf{secrezione ormonale}: secerne ossitocina e ormone
    antidiuretico;
  \item \textbf{effetti sul sistema nervoso autonomo}: è uno dei principali
    centri di integrazione del SNA; invia fibre discendenti ai nuclei del
    tronco encefalico che influenzano ritmo cardiaco, pressione arteriosa,
    diametro delle pupille, motilità gastrica, secrezione gastrica e altre
    funzioni;
  \item \textbf{termoregolazione}: i neuroni del nucleo preottico monitorano
    la temperatura corporea; quando questa si allontana dalla norma vengono
    attivati meccanismi per alzarla o abbassarla (sudorazione,
    vasodilatazione o vasocostrizione cutanea);
  \item \textbf{controllo dei ritmi circadiani}: la regione caudale
    ipotalamica fa parte della formazione reticolare, che contiene nuclei
    coinvolti nella regolazione del ritmo sonno-veglia;
  \item \textbf{assunzione di cibo e acqua}: i neuroni ipotalamici
    (osmocettori) controllano l'osmolarità del sangue e, in caso di
    disidratazione, innescano comportamenti finalizzati alla ricerca di
    acqua o cibo.
\end{enumerate}

% ---------------------------------------------------------------------------
\section{Il cervelletto}

Occupa la fossa cranica posteriore e rappresenta la seconda porzione più
estesa del cervello in toto: un ``encefalo mignon'', costituito da 2 emisferi
cerebellari connessi da uno stretto ponte chiamato \textbf{verme}.

\begin{itemize}
  \item presenta una corteccia esterna di sostanza grigia e una sostanza
    bianca interna, con morfologia simile a una felce, chiamata
    \textbf{arbor vitae};
  \item è collocato posteriormente al tronco dell'encefalo, a cui è collegato
    tramite i peduncoli cerebellari;
  \item forma il tetto del IV ventricolo.
\end{itemize}

\rubrica{Le tre regioni}
\begin{itemize}
  \item \textbf{corteccia}: sostanza grigia (più superficiale);
  \item \textbf{sostanza bianca} interna (più in profondità);
  \item \textbf{nuclei profondi}: materia grigia in profondità.
\end{itemize}

Contiene molte cellule di Purkinje, dalle complesse ramificazioni.

\rubrica{Dimensioni}
Forma grossolanamente ovoidale, con grande asse trasversale di 9\,cm,
diametro antero-posteriore di 6\,cm e verticale di 5\,cm; peso medio 140\,g.
È suddiviso nei lobi flocculonodulare, posteriore e anteriore.

\rubrica{Le cellule del Purkinje}
\begin{itemize}
  \item grossi neuroni il cui corpo oblungo è situato nell'omonimo strato;
  \item disposte a intervalli piuttosto regolari;
  \item possiedono un albero dendritico fittamente arborizzato, che origina
    da uno o due dendriti primari e si porta fino allo strato molecolare,
    dove contrae migliaia di sinapsi tramite numerose spine dendritiche con
    le fibre parallele delle cellule dei granuli;
  \item l'albero dendritico è caratteristicamente limitato a una sottile
    lamina di tessuto perpendicolare all'asse maggiore della lamella.
\end{itemize}

\rubrica{Ruolo funzionale}
Attaccato alla parte posteriore dell'encefalo, simile a un cavolfiore e
grande come un pugno; caratterizzato da altissima densità cellulare, è il
10\% del volume cerebrale ma contiene più o meno metà dei neuroni del SNC.
È coinvolto in \textit{movimento}, \textit{equilibrio} e \textit{postura};
oggi si ritiene che partecipi anche a memoria, umore, linguaggio, attenzione
e funzioni sensitive/emotive.

% ---------------------------------------------------------------------------
\section{Il midollo spinale}

Porzione extracranica del SNC, collocata all'interno del canale vertebrale.

\begin{itemize}
  \item comincia dal foro occipitale, come prosecuzione del bulbo, e continua
    in senso caudale fino al cono midollare a livello lombare;
  \item la sua estensione fibrosa, detta \textbf{filum terminale}, si
    prolunga fino al coccige;
  \item lunghezza media nell'adulto: 46\,cm (larghezza 1--1,5\,cm), in
    stretto rapporto con la lunghezza del tronco e variabile secondo sesso ed
    etnia;
  \item il livello di terminazione è variabile: nella maggior parte dei casi
    l'apice inferiore si trova nel disco intervertebrale tra L1 e L2, a
    livello del piano transpilorico, ma può terminare anche tra L2 e L3, o
    più raramente tra T12 e L1.
\end{itemize}

\rubrica{Organizzazione}
È formato da sostanza bianca esterna (prolungamenti dei neuroni) e sostanza
grigia interna (corpi dei neuroni), ed è diviso in quattro regioni --
cervicale, toracica, lombare e sacrale -- ciascuna formata da più segmenti.

\rubrica{Meningi spinali}
La dura madre riveste il midollo spinale formando un sacchetto di protezione,
irrorato al suo interno dal liquido cerebrospinale.

\rubrica{Legamenti denticolati}
Coppie di estensioni della pia madre spinale che collegano pia madre e
aracnoide spinali alla dura madre; impediscono i movimenti di scivolamento
verso il basso del midollo spinale.

\rubrica{Sostanza grigia}
A forma di H:
\begin{itemize}
  \item \textbf{commissura grigia}: contiene il canale centrale;
  \item \textbf{corna anteriori}: contengono i corpi cellulari dei
    motoneuroni;
  \item \textbf{corna posteriori}: consistono di interneuroni;
  \item la sostanza grigia è suddivisa tra regioni viscerali e somatiche.
\end{itemize}

\rubrica{Sostanza bianca}
Formata da assoni mielinici e amielinici, organizzati in tre tipi di fibre:
ascendenti, discendenti e commissurali.
